\section{评测与风险应对}
\begin{frame}{评测方法与主要风险}
\small
\begin{block}{评测计划}
对每个场景记录人工标准答案、Agent 结论、操作数、耗时、Token 与证据。判定一致率以全部有效评测场景为分母，无法判定单独统计且不算正确；对照实验固定应用版本、场景与预算，并进行重复运行。
\end{block}
\begin{tabularx}{\textwidth}{lX}
\toprule
风险 & 应对与验收关注点 \\
\midrule
定位失败或循环操作 & 限制恢复次数与整体运行时间，保留失败前观察，分析高频失败动作。\\
模型误判检查点 & 独立证据审查，对照人工标准答案，核查跨页证据与误报。\\
环境或数据状态漂移 & 固定应用版本与前置数据，记录重置步骤，验证新成员可复现。\\
场景、模型或进度受限 & 优先打通代表场景；正式评测前确认教师场景与模型映射，按周调整任务范围。\\
\bottomrule
\end{tabularx}
\end{frame}
